\chapter{Introduction} \label{chap:introchap} Historical analysis of atmospheric chemistry shows that global concentrations of ozone (O$_3$) are heavily influenced by multiple catalytic cycles involving halogen, nitrogen, hydrogen, and oxygen species \parencite{caffrey_cloud_2001, dobson_measurements_1923, he_sensitivity_1999}. The global movement of carbon among Earth’s atmosphere, biosphere, oceans, and geosphere, indirectly yet significantly influences tropospheric O$_3$ concentrations through greenhouse gas emissions and atmospheric chemical composition \parencite{danyang_ma_effect_2023,hata_impact_2023,liu_carbonyl_2022}. Ongoing regulation, monitoring, and development of air pollution control initiatives are crucial for maintaining safe exposure levels to O$_3$ \parencite{branis_long_2008,manes_regulating_2016,american_lung_association_disparities_2023,epa_benefits_2015}, as these levels can influence the evolutionary potential of biological mechanisms on Earth \parencite{grenfell_co-evolution_2010,kasting_life_2002}. This highly reactive, unstable, triatomic molecule; providing a canvas for dioxide cycles necessary for organic systems \parencite{caffrey_cloud_2001,lyons_rise_2014}, allows for redox reactions forming as both a product and constituent of the natural environment \parencite{chapleski_heterogeneous_2016,xing_representing_2016,zoran_ground_2014}. 

Anthropogenic sources have deeply affected these necessary cycles by altering natural concentrations of O$_3$ in unnatural ways \parencite{chauhan_rising_2023,cheadle_surface_2017,danyang_ma_effect_2023,flynn_relationship_2014}. As climate change continues to abet abnormal temperatures and weather patterns, the development of high-resolution predictive models, such as those presented in this thesis, becomes increasingly urgent to address the rapid rate of urbanization \parencite{balk_understanding_2018,epa_climate_2021,iglesias_risky_2021}. This thesis aims to provide a novel methodology that more effectively incorporates geospatial relationships into contemporary statistical approaches through residual kriging (RK) applied to three of the most populous counties in Arizona. The results of a chosen statistical method and residual krige (SMaRK) methodology are overlaid with statistical representations of a complex system, ozone O$_3$. Similar methods have already demonstrated great promise in modeling PM$_2.5$ for the contiguous United States \parencite{liu_improve_2018}. 
%%%%%%%%%%%%%%%%%%%%%%%%%%%%%%%%%%%%%%%%%%%%%%%%%%%%%%%%%%%%%%%%%%%%%%%%%%%%%%%%%%%%%%%%%%%%%%%%%%%%%%%%%%%%%%%%%%%%
\section{Machine Learning/Artificial Intelligence Predictions Require Incorporation of Geospatial Data} \label{introchap:mlneedgeo} Policies set by world governments can help reduce overall surface O$_3$ pollution, but their effectiveness depends on the models used to assess vulnerable areas. Current numerical modeling approaches, such as WRF-Chem, GEOS-Chem, GODfit Algorithms, and MUSICA \parencite{flandorfer_evaluation_2019,lerot_godfit_2010,lin_spatiotemporal_2016,yu_errors_2018}, are computationally expensive, which restricts the output of surface O$_3$ maps to coarse spatial resolutions unsuitable for public health studies. Health experts generate exposure data sets that provide highly accurate, non-rasterized depictions of surface O$_3$ \parencite{he_surface_2023,malley_updated_2017,turner_long-term_2016} by accounting for individual exposures. Most O$_3$ representations of the surface require numerous days to generate their final products \parencite{kwok_photochemical_2015,lin_spatiotemporal_2016,travis_systematic_2019,wu_apfoam_2021} and modeling efforts based on remote sensing measurements still struggle to capture the complexity of surface O$_3$ in urban settings without further improvements made by the researcher \parencite{balamurugan_tropospheric_2021,wang_evolution_2023}. When predictions from modeled O$_3$ reactions are compared to surface observations from monitors, they often struggle to simulate the variance among O$_3$ concentrations at resolutions below 500 m \parencite{alvarez-mendoza_spatial_2019,kasparoglu_spatial_2018}. This indicates that the underlying geospatial patterns used during observations may not fully characterize regions accurately, or that the resolution of the models is insufficient to simulate the photochemical environment of these finer observations with gathered imagery.

Machine learning methods, such as random forests, can capture complex relationships between O$_3$ and other covariates when generating rasterized predictions. However, they often fail to account for spatial correlations in remote sensing \parencite{liu_improve_2018,runmei_ma_full-coverage_2021,xiong_improving_2023}. With large, satellite-based datasets, linear and weighted regression tend to be exponentially faster for aggregating to higher spatial resolutions, at the cost of egregious error \parencite{feng_investigation_2023,huang_development_2017}. Combining an advanced model with its corresponding geospatial residuals helps account for the spatial correlation between modeled values while simultaneously accounting for complex relationships \parencite{liu_improve_2018}. RK is a commonly used geostatistical method based on multiple linear regressions and variograms \parencite{meng_assessment_2013}.

In application to surface O$_3$ values, this method helps reduce bias from covariates while incorporating a geospatial perspective that acknowledges spatial heterogeneity among these features. Five ensemble learning models are compared: Adaptive boosting, Gradient Boosting, Extreme Gradient Boosting, Random Forest, and Multi-Layered Perceptron. These were created over a study area containing Arizona’s most populous counties of Maricopa, Pinal, and Pima from 2019-2024. Each model was created using a hyper-parameter tuner available through sci-kit learn in python \cite{buitinck_api_2013}. 
%%%%%%%%%%%%%%%%%%%%%%%%%%%%%%%%%%%%%%%%%%%%%%%%%%%%%%%%%%%%%%%%%%%%%%%%%%%%%%%%%%%%%%%%%%%%%%%%%%%%%%%%%%%%%%%%%%%%
\section{Why Surface O$_3$? – Relations to Urban Air Quality} \label{introchap:whyO$_3$} O$_3$ has been detected in Earth’s stratosphere, troposphere, and more recently, on its surface \cite{claeyman_thermal_2011,davies_episodes_1994,he_sensitivity_1999,lin_springtime_2012,gao_case_2016}. In 1955, the United States issued the Air Quality Control Act which dedicated the nation’s vast resources to atmospheric monitoring, addressing air pollution concerns post-scientific revolution enticed by the Space Race of 1950 \cite{epa_benefits_2015}. The acting President of the United States at the time; Richard Nixon, aided in the establishment of an Environmental Protection Agency (EPA) and National Oceanic and Atmospheric Administration (NOAA) to combat this newfound enemy with extreme prejudice \cite{nixon_reorganization_1970}. As a result, new environmental movements have emerged due to the unprecedented economic growth of that period; regulatory policies that initially hindered some sectors later demonstrated long-term beneficial health and economic trends, as theorized in the early 1990s by Dr. Michael E. Porter \parencite{ambec_theoretical_2002}. Geographic Informational Science (GIScience) has been dedicated to the proper implementation of these goals set forth in the early 1970s. With it came a new way to write Earth’s history via digital transformations of newfound resources; relating to trends near and far to answer difficult questions associated with life. From micro- to macro-scales, a geographer specializes in observations and predictions from planetary to microbial systems, inadvertently weaving complex knowledge of space and time \cite{goodchild_geographical_1992}. 

Many observations have been made with the wealth of material garnered from this time period regarding communities at risk due to high surface O$_3$ concentrations. This reaction is exacerbated by urban heat island effects, vehicle emissions, and industrial processes that contribute to a variety of pollutants resulting from the mechanisms that drive it \parencite{afonso_characterization_2017,li_effects_2014}. It is a precursor to many things due to its highly degenerative state \parencite{afonso_characterization_2017}. When large amounts of UV radiation interact with these sources, O$_3$ begins to form and can have numerous effects on organic and in-organic material exposed to it \parencite{chapleski_heterogeneous_2016,huang_biogenic_2024,brauer_global_2024}. Coastal cities typically demonstrate stronger associations with natural vegetation by providing non-spatial heterogeneity when compared to anthropogenic sources, reflecting conditions prior to human development \parencite{cai_oxidative_2022};\cite{chen_analysis_2024,martins_spatial_2016}. In contrast, Phoenix and Tucson (PHOTUC) are characterized by arid-dry climates that exhibit spatial heterogeneity in constituents associated with human development, influenced by dry deposition from local farming processes and emissions \parencite{chiacchiaretta_inland_2024,ding_contrasting_2023,le_tourneau_socio-ecological_2024,ye_diagnosing_2022}. 

The area of interest (AOI) encompassing PHOTUC presents significant challenges in accurately representing O$_3$ due to numerous anthropogenic factors masking non-heterogeneity in among covariates. The predictive performance of geostatistically enhanced numerical models is analyzed in this area due to the challenge advancements to computational systems. Global representations of surface O$_3$ are offered in this era of Big Data \cite{cao_deep_2022,bughin_big_2016,curry_big_2016,xu_mapping_2023}. This thesis applies the three laws of Geography to ultimately propose a high-resolution representation of surface O$_3$ at tiled resolutions ranging from 300m down to 50m for use in rural to urban ecologies respectively. Surface O$_3$ has been found to be one of the leading pollutants on the Global Burden of Disease (GBD) \parencite{brauer_global_2024,liu_toward_2024,sang2022}, harboring stealthy consequences that can only be seen with appropriate consideration of both formation and exposure to fully understand its impacts \parencite{watson_machine_2019,abdullah_relationship_2017,carvalho_o3_2022}. 

Varying concentrations of O$_3$ have been found to have numerous beneficial and adverse side-effects on oxide-systems in a myriad of ways due to interactions with surrounding ecosystems \parencite{ding_investigating_2021,epa_integrated_2020,manes_regulating_2016}. Upon closer examination, ground-level O$_3$ has been shown to have a lasting impact on human health, manifesting as reductions in life expectancy \parencite{zhang_effect_2019,turner_long-term_2016}, abrupt changes to normalized Earthly air chemistry \parencite{lin_spatiotemporal_2016,wu_hybrid_2023}, and alterations in related oxide-based cycles \parencite{barzeghar_long-term_2020,ding_investigating_2021}. Additionally, surface O$_3$ exposure contributes to several adverse health outcomes, including respiratory symptoms, childhood cancers, adverse birth outcomes, overall mortality, and neurological disorders \parencite{buadong_association_2009,ghozikali_effect_2015,tang_impact_2024,turner_long-term_2016}. These studies and many more across the world( \cite{brown-steiner_asian_2011,cheng_estimator_2018,girach_changes_2012} have also found the most significant constituents of ground-level O$_3$ reactions to be nitrogen-oxides (NOx) and volatile organic compounds (VoCs). Many of these interactions stem from anthropogenic sources \parencite{danyang_ma_effect_2023,chauhan_rising_2023} in-lieu of natural interactions found before settlement occurred in these areas. 
%%%%%%%%%%%%%%%%%%%%%%%%%%%%%%%%%%%%%%%%%%%%%%%%%%%%%%%%%%%%%%%%%%%%%%%%%%%%%%%%%%%%%%%%%%%%%%%%%%%%%%%%%%%%%%%%%%%%
\section{A Tale of Two Layers} \label{introchap:twolayers} O$_3$ is highly reactive and decomposes rapidly, increasing the likelihood of reactive oxygen species (ROS) existing in any environment where it is present \parencite{voter_ozone_2001,bocci_studies_1998}. It is generally beneficial to the natural environment at consistent levels, serving as a semi-catalyst for certain biological and chemical systems \parencite{wiley_ozone_2017}. The patterns observed by researchers regarding surface O$_3$ (i.e., 2m above the ground) have been challenging to trace due to the numerous interactions it has with these systems. Highly concentrated areas of O$_3$ near Earth’s surface typically impact populations outside of highly industrialized zones characterized by dense vegetation, as the byproduct of plants is mainly carbon dioxide (CO$_2$) \parencite{lei_impacts_2024,sadiq_effects_2016};\cite{yadong_lei_indirect_2021}. Studies interested in surface O$_3$ exposure within one’s activity space have noted severe differences in health outcomes among varying populations \parencite{crouse_ambient_2015,ekland_effect_2021,malley_updated_2017,fu_prenatal_2020,tang_impact_2024}. To better understand the harmful cycles of surface O$_3$, the underlying principles determining its formation in both layers of Earth’s atmosphere are mentioned. 

Stratospheric O$_3$ protects the surface of Earth by absorbing UV rays, maintaining O$_3$ levels seen on the surface \parencite{dobson_measurements_1923,hodzic_response_2018}. Due interactions with chemical compositions in this layer, O$_3$ was found to naturally follow seasonal cycles which correlate to positioning of Earth’s tilt in its rotational pattern and distance to the sun. However, due to rising temperatures from anthropogenic activities, this beneficial, and natural cycle is threatened as it begins mixing with surface air quality and chemistry \parencite{akritidis_evaluating_2013,cristina_penache_temporal_2019,sadiq_effects_2016}. As the O$_3$ Layer doesn’t absorb all UV radiation, what’s left over reflects off the surface of the Earth and allows for the reaction to take place. 

Ground-level O$_3$ combined with other reactive gases will tend to correlate to the quality of air encompassing human communities \parencite{gaudel_tropospheric_2018};\cite{schultz_tropospheric_2017}. As this dual formation involves titrations with varying chemicals, complex, non-linear associations begin closer to the source of common pollutants; later, they are reincorporated into the environment at some distance away. This makes it difficult for these O$_3$ reactions to be identifiable and effectively modeled for decision making and regulatory entities \parencite{abdullah_development_2019,balamurugan_tropospheric_2021,huang_development_2017}. Furthermore, stratospheric O$_3$ concentrations in the future will depend on the decrease of O$_3$ depleting substances (ODS) found near the surface due to the reduction of the frequency and hence; likelihood of these interactions \parencite{huang_development_2017,mader_statistical_2007}. 

While most hazardous pollutants are emitted directly, complex pollutants such as surface O$_3$ form as byproduct, later becoming a constituent for other chemicals. O$_3$ is known as a secondary pollutant \parencite{lin_modeling_2012,watson_machine_2019,venkanna_environmental_2015} wherein it can act as an ingredient and result of chemical processes. The slightest change in wind speeds, environmental conditions (coastal to arid), and ecological elements (available greenspace vs. concrete) can influence O$_3$ formation and degradation \parencite{badia_modelling_2023,meng_influence_2022,xing_representing_2016}. The data provided by these, and numerous other works gives insights into historical transport of O$_3$ due to urban, suburban, and rural development which were incorporated during development.
%%%%%%%%%%%%%%%%%%%%%%%%%%%%%%%%%%%%%%%%%%%%%%%%%%%%%%%%%%%%%%%%%%%%%%%%%%%%%%%%%%%%%%%%%%%%%%%%%%%%%%%%%%%%%%%%%%%%
\section{Surface O$_3$ Exposure and Transport} \label{introchap:surfeat} Studies which gathered similar data products have found several links between urbanization and populations which are exposed to unhealthy levels of air pollution above those of current standards set by the Environmental Protection Agency and World Health Organization \parencite{epa_climate_2021,kumar_integrated_2015,liu_impact_2022}. The EPA establishes an AQI for communicating daily air quality in which other major air pollutants are also regulated under the Clean Air Act. O$_3$ has direct interactions with all pollutants labeled hazardous by these agencies due to its molecular instability. The trajectory of future emissions has a direct impact on policy decisions affecting numerous sociol-economic statuses (SES) due to inherent constituents. 

Sand, smoke, volcanic plumes, dust, other gaseous components like clouds can affect the overarching operations of O$_3$ cycles and complexity as seen in many studies \parencite{harithasree_surface_2024,venkanna_environmental_2015,tong_characteristics_2017}. The precursors for its existence are risks to human health, as indicated by informed policies limiting exposures to other pollutants (e.g. Particulate matter (PM), Carbon monoxide (CO), Nitrogen dioxide (NO$_2$), etc.) per the CDC, EPA, and WHO \cite{cdc_air_2024,epa_benefits_2015,liu_toward_2024}. Limiting the general reduction of emissions like CO$_2$, formaldehyde (HCHO or CH$_2$O), methane (CH$_4$) and nitrous oxides (NO$_x$), tend to reduce surface O$_3$ reactions by decreasing the probability for a reaction to occur in populated areas with large albedo \parencite{al-harbi_spatiotemporal_2020,badia_modelling_2023,wu_hybrid_2023}. 

Current representations and remote sensing methods used for surface O$_3$ can be too coarse for urban analysis, and might not highlight key details known to interact with urbanization and natural disasters in communities at finer spatial scales \parencite{nawaz_source_2023,zhang_effect_2019,abdullah_development_2019}. These instruments often rely on innovative approaches to aerosol forecasting and implement various corrective algorithms for accurate representation of meteorological variables. A Statistical Model and Residual Kriging methodology is proposed to model surface O$_3$ estimates with a spatial resolution of 250m, which is appropriate for urban analysis \parencite{wang_does_2023}. This thesis aims to compare the performance of common machine learning (ML) and artificial intelligence (AI) methods with the enhancements provided by the residual Kriging method on forecasted O$_3$ values for further development into a Python-based air pollution modeling library.
%%%%%%%%%%%%%%%%%%%%%%%%%%%%%%%%%%%%%%%%%%%%%%%%%%%%%%%%%%%%%%%%%%%%%%%%%%%%%%%%%%%%%%%%%%%%%%%%%%%%%%%%%%%%%%%%%%%%
\section{Structure of Thesis} \label{introchap:structure} Beginning with a comprehensive literature review on the formation of O$_3$, it is evident that numerous aspects can affect local ecologies, health, and communities. These are extensively discussed to provide a robust foundation for feature creation and model tuning based on scientific evidence. Some drivers included in the initial dataset were not utilized in the final model; however, all data initially gathered was derived from the synthesis of literature. A section dedicated to data sources and materials offers further justification for the features best suited for ML/AI ensembles and their integration with the RK method. 

After understanding the fundamental drivers, constituents, and numerical model tuning methods, a geospatial regression krige is applied to the resulting uncertainty left by the estimated trend. The full combination of these techniques creates daily high spatial resolution rasters for the urban cities of Phoenix and Tucson in Arizona. These cities are primarily encompassed in two counties, with a third situated between them. Extremely detailed demographic maps illustrate this along with income, total population, and occupied households, utilizing census data provided by the US Census Bureau to highlight the full potential of the final images. 

The main portion of the thesis concludes with a detailed discussion of future directions, model improvements, and further justification for the necessity of high-spatial resolution surface O$_3$ models. Additionally, it introduces a spatial program dedicated to rapidly assigning air pollution exposures. This program emerges as a product of this project, aimed at enhancing health-related studies conducted in the past five years, particularly within the contiguous United States. A brief section on the application of geographical laws and theories in historical scientific literature serves as the final remarks, illustrating a proposed economic solution derived from advancements in air pollution modeling. The methodology employed in this thesis is best characterized as a common case study that utilizes the scientific method, geospatial statistical analysis, and the application of the three main laws of geography to contemporary ML/AI modeling methodologies.